% Source-aligned English for printed p.186.
\clearpage
\SourcePage{186}

\LayerTarget{tohoku:proof:proposition-4.1.1:p1}
\LayerTarget{tohoku:diagram:4.1.1:finite-type-stalk-hom}
\[
\begin{array}{ccccc}
0&\to&\mathbf{Hom}_{\mathbf O}(A,B)(x)&\longrightarrow&
  \mathbf{Hom}_{\mathbf O}({\mathbf O}^n,B)(x)\\
&&\big\downarrow&&\big\downarrow\\[0.45em]
0&\to&\mathrm{Hom}_{{\mathbf O}(x)}(A(x),B(x))&\to&
  \mathrm{Hom}_{{\mathbf O}(x)}({\mathbf O}^n(x),B(x)).
\end{array}
\]
\LayerTarget{tohoku:eq:4.1.4}
Now $\mathbf{Hom}_{\mathbf O}({\mathbf O}^n,B)=\mathbf{Hom}_{\mathbf O}({\mathbf O},B)^n=B^n$, since there is a
canonical isomorphism
\[
\tag{4.1.4}
\mathbf{Hom}_{\mathbf O}({\mathbf O},B)=B.
\]
\LayerTarget{tohoku:proof:proposition-4.1.1:p2}
The groups in the last column of the diagram therefore identify respectively
with $B(x)^n$ and $\mathrm{Hom}_{{\mathbf O}(x)}({\mathbf O}(x),\mathrm B(x))^n=B(x)^n$. Thus the last
vertical homomorphism is an isomorphism, and it follows at once that the first
homomorphism is a monomorphism. Now suppose that $A$ is pseudo-coherent at
$x$; that is, after restricting if necessary to a suitable neighborhood of
$x$, there is an exact sequence
${\mathbf O}^m\to {\mathbf O}^n\to A\to0$. Again using the
left exactness of the $\mathrm{Hom}$ functors, we obtain a homomorphism of
exact sequences:
\LayerTarget{tohoku:diagram:4.1.1:pseudo-coherent-stalk-hom}
\[
\begin{array}{ccccccc}
0&\to&\mathbf{Hom}_{\mathbf O}(A,B)(x)&\longrightarrow&
  \mathbf{Hom}_{\mathbf O}({\mathbf O}^n,B)(x)&\longrightarrow&\mathbf{Hom}_{\mathbf O}({\mathbf O}^m,B)(x)\\
&&\big\downarrow&&\big\downarrow&&\big\downarrow\\[0.45em]
0&\to&\mathrm{Hom}_{{\mathbf O}(x)}(A(x),B(x))&\to&
  \mathrm{Hom}_{{\mathbf O}(x)}({\mathbf O}^n(x),B(x))&\to&
  \mathrm{Hom}_{{\mathbf O}(x)}({\mathbf O}^m(x),B(x)).
\end{array}
\]
\LayerTarget{tohoku:proof:proposition-4.1.1:p3}
We have already noted that the last two vertical homomorphisms are bijective.
It follows at once that the first vertical homomorphism is an isomorphism.

\LayerTarget{tohoku:question:injectivity-of-stalk}
Now suppose that $B$ is an injective object of the abelian category $\mathbf C^{\mathbf O}$.
Can we conclude that $B(x)$ is an injective ${\mathbf O}(x)$-module? By Lemma 1 of 1.10,
it is enough to verify that every ${\mathbf O}(x)$-homomorphism from a left ideal $\mathbf a$ of
${\mathbf O}(x)$ into $B(x)$ extends to an ${\mathbf O}(x)$-homomorphism from ${\mathbf O}(x)$ into $B(x)$.
Since the restriction of an injective ${\mathbf O}$-module to an open subset is still
injective (Prop. 3.1.3), it is enough to find an open neighborhood $U$ of $x$,
a submodule $A$ of ${\mathbf O}|U$, and a homomorphism $u$ from $A$ to $B|U$ such that
$A(x)=\mathbf a$ and the homomorphism from $A(x)$ to $B(x)$ defined by $u$ is the
given homomorphism $\mathbf a\to B(x)$. Moreover, if $A$ is
pseudo-coherent at $x$, the existence of $u$ follows automatically from
Proposition 4.1.1. We therefore prove:

\LayerTarget{tohoku:lemma:4.1.1}
\textsc{Lemma 1.} \textit{Suppose that ${\mathbf O}$ is a coherent sheaf of
rings \textup{[15, Chap. I, §2]}, that $M$ is a coherent ${\mathbf O}$-module \textup{(}a ``coherent sheaf
of modules'' in the terminology of \textup{[15])}, and that $\mathbf n$ is a submodule of
finite type of $M(x)$. Then, on some suitable open neighborhood $U$ of $x$,
there is a coherent subsheaf $N$ of $M$ such that $N(x)=\mathbf n$.}

\LayerTarget{tohoku:proof:lemma-4.1.1}
Let $(n_i)$ be a finite system of generators of $\mathbf n$ ($1\leqq i\leqq k$). For
each $i$, we have $n_i\in M(x)$, hence $n_i=f_i(x)$, where $f_i$ is a section
of $M$ defined on an open neighborhood of $x$. We may suppose that all the
$f_i$ are defined on the same open neighborhood $U$; over $U$ they then
define a homomorphism ${\mathbf O}^k\to M$. Take $N$ to be the image of
this homomorphism, which is coherent as a submodule of finite type of the
coherent sheaf $M$.

\TohokuContinuation{187}
